\SourceLeaf{414}%
Aus \(-\frac12b[2]+\frac12a[3]\) und \(\frac12B[2]-\frac12A[3]\) folgt, wenn man zur Abkürzung
\[\begin{gathered}aB-bA=\lambda,\quad ab(A+B)=f,\quad\tfrac12AB(A+B)+\lambda n=g\\
AB(a+b)=F,\quad\tfrac12aBB+\tfrac12bAA=G\end{gathered}\]
schreibt,
\[\begin{gathered}\lambda T+G=\lambda(t+f)x,\quad\lambda t-g=(\lambda T-F)x\quad\text{oder}\quad\lambda(t-Tx)=g-Fx\end{gathered}\]
und damit
\[\begin{aligned}\lambda t&=\frac{+g-(F+G)x+fxx}{1-xx}\quad.\ .\ .\ .\ .\ .\ .\ .\ [4]\\
\lambda T&=\frac{-G+(f+g)x-Fxx}{1-xx}\quad.\ .\ .\ .\ .\ .\ .\ .\ [5]\end{aligned}\]
\par
Gleichung \([1]\) in der Form \((t-Tx)^2+TT(1-xx)=rr\) ergibt
\[\lambda\lambda rr-(g-Fx)^2=\lambda\lambda TT(1-xx)\quad.\ .\ .\ .\ .\ .\ [6]\]
\par
Setzt man darin den Wert von \(\lambda T\) aus \([5]\) ein, erhält man die kubische Gleichung
\[\begin{aligned}&2fFx^3-(ff+2fg+FF+2FG+\lambda\lambda rr)xx\\
&\quad+(2fG+2gF+2gG)x+\lambda\lambda rr-gg-GG=0\quad.\ .\ [7]\end{aligned}\]
\par
Nachdem \(x\) auf diese Weise bestimmt ist, findet man die Koordinaten des Punktes \((2)\) aus den obigen Ausdrücken. Setzt man dort für \(t-Tx\) und \(T\) die zuvor angegebenen Werte ein, nehmen sie folgende Form an
\[n-\frac{g-Fx}{\lambda},\qquad\frac{-G+(f+g)x-Fxx}{\lambda\BareRoot{(1-xx)}}\]
und man findet die beiden anderen Punkte \((1)\) und \((3)\), indem man zu diesen Werten jeweils \(-Ax,-A\BareRoot{(1-xx)}\) und \(+Bx,+B\BareRoot{(1-xx)}\) addiert.

Da jedem Kosinus zwei Winkelwerte entsprechen, oder, was dasselbe ist, da der Wurzelausdruck \(\BareRoot{(1-xx)}\) sowohl positiv als auch negativ genommen werden kann, liefert jede zulässige Wurzel der Gleichung zwei Lösungen, nämlich zwei zur Geraden \(\mathrm{(II)}\) symmetrische Lagen der Punkte \((1),(2),(3)\). Das ist schon für sich genommen klar. Ist allerdings \(+1\) oder \(-1\) eine Wurzel der Gleichung \([7]\), lässt sich die obige Formel für die Ordinate nicht verwenden, weil Zähler und Nenner dann null werden. Stattdessen muss man
